% ============================================================
% corex-preamble.tex
% Preamble bersama untuk semua dokumen formal Corex.
% Di-\input dari tiap dokumen (Proposal, NDA, Perjanjian Kerja, dst)
% supaya perubahan styling cukup dilakukan di SATU tempat.
% ============================================================

\pdfminorversion=7

\usepackage[a4paper,margin=2.5cm,top=3.2cm,bottom=3cm,headheight=22pt]{geometry}
\usepackage{xcolor}
\usepackage{graphicx}
\usepackage{fancyhdr}
\usepackage{titlesec}
\usepackage{enumitem}
\usepackage{array}
\usepackage{longtable}
\usepackage{booktabs}
\usepackage{hyperref}
\usepackage{tcolorbox}
\usepackage{parskip}
\usepackage{fontenc}
\usepackage[T1]{fontenc}
\usepackage{lmodern}
\usepackage{libertine}
\usepackage{lastpage}

% ---------- Palet Warna Corex ----------
% Navy: warna dominan teks judul & garis. Signal Cyan: aksen SEMPIT saja
% (heading rule, link, badge) -- konsisten dengan aturan di DESIGN-DNA.md,
% supaya dokumen resmi pun tidak jadi "warna Corex dominan di semua tempat".
\definecolor{corexnavy}{HTML}{0B0F19}
\definecolor{corexsignal}{HTML}{00B39E}  % sedikit digelapkan dari #00D9C0 supaya kontras cukup di atas kertas putih
\definecolor{corexgray}{HTML}{5B6472}
\definecolor{corexlinegray}{HTML}{D5D9E0}

% ---------- Judul Section ----------
\titleformat{\section}
  {\normalfont\Large\bfseries\color{corexnavy}}
  {\thesection}{0.6em}{}
\titleformat{\subsection}
  {\normalfont\large\bfseries\color{corexnavy}}
  {\thesubsection}{0.6em}{}
\titlespacing*{\section}{0pt}{1.4em}{0.8em}
\titlespacing*{\subsection}{0pt}{1.1em}{0.6em}

% ---------- Hyperlink tanpa kotak norak ----------
\hypersetup{
  colorlinks=true,
  linkcolor=corexnavy,
  urlcolor=corexsignal,
  citecolor=corexnavy,
  pdfcreator={Corex},
}

% ---------- Komando Kustom ----------
% Kotak disclaimer wajib (dipakai corex-legal di NDA/Perjanjian Kerja)
\newtcolorbox{corexdisclaimer}{
  colback=corexnavy!4,
  colframe=corexnavy,
  boxrule=0.8pt,
  arc=1pt,
  left=8pt, right=8pt, top=6pt, bottom=6pt,
  fontupper=\small
}

% Placeholder yang wajib diisi manusia -- dibuat MENCOLOK sengaja,
% supaya draft tidak sengaja terkirim ke klien dengan bagian kosong.
\newcommand{\CorexFill}[1]{\textcolor{red!70!black}{\textbf{[#1]}}}

% Badge kecil status dokumen (DRAFT / FINAL) di title page
\newcommand{\CorexStatusBadge}[1]{%
  \colorbox{corexnavy}{\textcolor{white}{\small\bfseries\ #1\ }}%
}

% Blok tanda tangan standar
\newcommand{\CorexSignatureBlock}[2]{%
  \begin{minipage}[t]{0.45\textwidth}
    \vspace{2.5cm}
    \hrule
    \vspace{4pt}
    \textbf{#1} \\[2pt]
    #2
  \end{minipage}
}
